\begin{abstract}
We prove that local transport on the flag space of the combinatorial
dodecahedron determines a connected $60$-vertex tetravalent graph
$G_{60}$, canonically identified with the arc-transitive graph
$\mathrm{AT4val}[60,6]$ (House of Graphs: Graph52002).

The construction admits natural $2$-fold quotients
\[
G_{60} \to G_{30} \to G_{15},
\]
where $G_{15} \cong L(\mathrm{Petersen})$.

We show that the lifted transport induces a canonical
$\mathbb{Z}_2$-valued cocycle on $G_{15}$, called the
\emph{Thalean cocycle}, which governs the signed $2$-lift
$G_{30} \to G_{15}$ and represents a nontrivial cohomology class.

From the transport system we construct a sector--edge incidence operator
$M \in \mathbb{R}^{15 \times 30}$ satisfying the polynomial identity
\[
M M^{\mathsf T} = A^3 + 2A^2 + 2I,
\]
where $A$ is the adjacency matrix of $G_{15}$.

After centering and normalization, this yields a spherical embedding of
$G_{15}$ in $\mathbb{R}^{14}$ with exactly three distinct inner products,
determined by graph distance.

Thus the transport system decomposes into two complementary structures:
a discrete $\mathbb{Z}_2$-valued cocycle capturing sheet holonomy, and a
rigid spherical sector geometry arising from global incidence. Together
these encode both the obstruction and the metric structure of the
Petersen line-graph core.
\end{abstract}
